\chapter{系统详细设计与实现}

\section{开发环境与工具}

本系统采用前后端分离方式进行开发与部署。后端以 Java 17 和 Spring Boot 2.7 为核心开发框架，借助 Maven 完成依赖管理和项目构建；前端采用 Vue 3 与 Vite 构建单页应用；数据层使用 MySQL 管理业务数据；开发过程中通过浏览器开发者工具和接口调试工具对请求响应与页面交互进行辅助验证。系统开发环境如表~\ref{tab:dev-env}所示。

\begin{table}[htbp]
\centering
\footnotesize
\caption{系统开发环境与工具}
\label{tab:dev-env}
\begin{tabular}{|c|c|c|}
\hline
类别 & 名称 & 说明 \\\hline
操作系统 & Windows & 系统开发与论文撰写环境 \\\hline
后端语言 & Java 17 & 后端业务逻辑开发 \\\hline
后端框架 & Spring Boot 2.7 & RESTful 服务构建与模块整合 \\\hline
安全框架 & Spring Security & 认证授权与访问控制 \\\hline
持久化技术 & Spring Data JPA & 实体映射与数据库访问 \\\hline
数据库 & MySQL & 业务数据持久化存储 \\\hline
前端框架 & Vue 3 & 页面构建与交互实现 \\\hline
构建工具 & Vite、Maven & 前后端工程构建与依赖管理 \\\hline
通信技术 & SockJS、STOMP & 实时消息推送与订阅 \\\hline
HTTP 调用 & Axios & 前后端接口访问封装 \\\hline
\end{tabular}
\end{table}

\section{各功能模块实现}

\subsection{用户认证与账号管理模块实现}

用户认证与账号管理模块是系统运行的入口。后端通过 Spring Security 构建安全过滤链，在登录请求进入系统后，由认证管理器完成用户名和密码校验，认证成功后生成 JWT 返回前端。前端在登录成功后保存令牌，并在 Axios 请求拦截器中统一附加到后续请求头，实现多页面访问下的身份延续。后端在每次请求到达时解析令牌并恢复用户身份信息，将角色集合注入安全上下文，从而支撑后续接口鉴权。

账号管理功能主要面向系统管理员。管理员可以维护用户基本信息、角色归属和账号状态。为防止越权，涉及账号新增、停用和角色变更等操作均在后端服务层进行角色校验，并对关键更新操作记录审计信息。前端页面则通过角色控制决定功能入口是否可见，实现页面限制与接口限制双重防护。

\subsection{顾客下单与订单管理模块实现}

顾客下单模块以前端订单表单为载体，支持产品选择、数量填写、订单提交和历史订单查询。前端在提交前对订单项完整性进行基础校验，然后通过 Axios 向后端订单接口发送请求。后端控制层接收订单数据后，调用服务层完成订单主表与订单明细表写入，并初始化订单状态为待审核。

订单管理模块贯穿系统多个角色。顾客可查看订单处理进度，销售管理员可执行审核，仓库管理员可基于订单触发出库或生产联动。为保证业务一致性，订单状态不允许任意跳转，而必须按照既定流程由服务层统一推进。这样能够避免前端直接修改状态造成的链路断裂问题。

\subsection{销售管理模块实现}

销售管理模块的核心在于订单审核与销售记录维护。销售管理员进入审核页面后，可查看待审核订单详情，对订单有效性、数量合理性和交付约束进行确认。审核操作提交后，后端服务层根据结果更新订单状态，并生成对应的销售处理记录。当审核通过时，系统进一步向仓库管理员推送待处理通知；当审核不通过时，则将结果反馈给顾客。

从实现机制看，销售管理模块通过“订单主表状态更新 + 销售处理记录写入 + 消息通知推送”的组合方式完成业务闭环，既保留了审核留痕，也为下游角色提供了明确的处理依据。

\subsection{仓库管理与库存预警模块实现}

仓库管理模块包含库存查询、入库、出库、发货和库存预警等功能。系统在数据库中维护成品和原材料库存实体，并通过库存流水表记录每次数量变化的原因、数量和操作人员。仓库管理员在执行入库、出库或发货操作时，前端向后端提交库存变更请求，后端在事务中同步更新当前库存数量与流水记录，以保证结果一致。

库存预警模块通过定期计算或业务触发检查库存状态。当当前数量叠加在途数量后仍低于安全库存阈值时，系统生成预警信息并写入消息表，同时通过实时通知机制推送给仓库管理员和采购管理员。该模块的实现重点在于将库存风险识别从人工经验判断转变为规则化处理，为生产和采购协同提供前置信号。

\subsection{生产管理模块实现}

生产管理模块负责承接库存不足或周计划下发后的生产需求。后端服务根据订单缺口或计划结果生成生产计划记录，并分配给生产管理员。生产管理员可在前端页面中查看待执行任务、填报执行进度和维护生产结果。生产结束后，系统将对应任务状态更新为待质检，并触发质检通知。

在实现层面，生产管理模块的关键在于对任务状态进行严格控制。生产计划、生产执行、生产完成和待质检等状态由服务层统一推进，并与订单履约需求和库存状态联动，从而保证生产活动不是孤立存在，而是服务于订单交付与库存补充。

\subsection{采购与供应商管理模块实现}

采购与供应商管理模块支持采购申请、采购订单生成、供应商接单、发货反馈和仓库收货入库等业务。采购管理员在前端页面中可查看由库存预警或周计划触发的采购建议，也可手动发起采购申请。申请经确认后由后端生成采购订单，并将订单推送给对应供应商。供应商登录后可查看待处理订单，执行接单和发货确认操作。仓库管理员在货物到达后完成收货入库，必要时可与质检模块联动进行来料检验。

该模块的实现重点在于企业内部与外部供应商之间的信息贯通。系统不仅记录采购单据本身，还通过状态更新与消息通知保持采购管理员、供应商和仓库管理员之间的信息同步，从而降低人工跟踪成本。

\subsection{质检与质量追溯模块实现}

质检模块主要用于处理生产完成后的成品质检以及必要的来料质检。质检管理员在系统中查看待检任务，录入检验结果、检验时间和备注信息。若检验合格，系统允许对应产品入库或进入后续发货流程；若检验不合格，则阻断后续流转并保留异常记录。

质量追溯模块则基于订单号或关联业务标识整合订单信息、生产记录、质检记录和库存流转记录。顾客在追溯页面输入订单信息后，后端按链路查询相关记录并返回结果。通过该机制，系统能够对关键业务环节形成可查询、可核验的数据链条，提高交付透明度。

\subsection{实时消息通知模块实现}

实时消息通知模块基于 WebSocket/STOMP 构建。后端在订单审核通过、库存预警触发、采购订单下发、供应商发货、生产任务生成、质检完成等业务节点，向指定主题或用户消息通道发送通知。前端在用户登录后建立 WebSocket 连接，并订阅与当前角色相关的消息主题。消息到达后，前端可刷新待办数量、展示通知列表或更新页面局部状态。

为了保证消息可追溯，系统不仅进行即时推送，还将通知内容写入消息通知表。这样即使用户在离线状态下未实时接收消息，也可在下次登录时查看历史通知，从而兼顾实时性与留痕性。

\subsection{周计划智能生成模块实现}

周计划生成模块并非依赖复杂智能算法，而是基于业务规则对未来一周的供需关系进行计算。后端服务首先汇总订单需求、当前库存、在途数量、安全库存和历史消耗数据，再结合产品与原材料之间的 BOM 关系估算未来一周的生产量和采购量建议。系统生成结果后，生产管理员和采购管理员可在对应页面中查看并作为业务决策参考。

在实现链路上，该模块通过计划服务统一读取订单、库存、采购和产品配置数据，完成规则计算后将结果写入生产周计划表与采购周计划表，并向相关角色发送通知。该设计体现了系统从事务处理向辅助决策延伸的工程价值。